\bookchapter{Penutup: Kekal Ingin Tahu, Kekal Berwaspada}

Tahniah, anda telah membaca sehingga sini.

Buku ini dari perkakasan telefon hingga ekosistem perisian, dari tetapan harian hingga mod pembangun, dari perlindungan keselamatan hingga harta digital. Jika ia bantu selesaikan satu masalah praktikal - walaupun hanya belajar tutup satu kebenaran yang tidak patut dibuka, membentuk satu tabiat sandaran - maka perkataan ini tidak sia-sia.

Tetapi lebih penting daripada apa-apa kemahiran khusus, adalah dua perkara:

\begin{itemize}
  \item \textbf{Kekal ingin tahu}: Telefon berubah setiap tahun, AI berubah, skrin lipat berubah, komunikasi satelit berubah. Langkah operasi yang dipelajari hari ini mungkin esok sudah ketinggalan, tetapi rasa ingin tahu "mahu tahu kenapa" tidak akan pernah ketinggalan.
  \item \textbf{Kekal sedar}: Telefon adalah alat, bukan tuan. Ia direka untuk menarik perhatian anda, menguasai masa anda, membimbing perbelanjaan anda. Guna ia adalah pilihan anda, diguna olehnya adalah hilang kawalan. Secara berkala pandang masa skrin, secara berkala fikir "saya sedang guna telefon buat sesuatu, atau sedang dibuang masa oleh telefon".
\end{itemize}

Teknologi berubah, tetapi keperluan manusia tidak berubah - kita perlukan rasa selamat, perlukan sambungan, perlukan dihormati. Telefon hanya satu pembawa. “Kelas telefon” sebenar tidak ada dalam buku, pada setiap saat anda buka kunci skrin.

\begin{tcbpractice}{Latihan terakhir}
Tutup buku ini, ambil telefon anda, buat satu perkara:\textbf{Guna 30 saat fikir, telefon anda ada tiga perkara apa yang “kehilangan akan betul-betul sakit hati”?}Kemudian, pastikan ia telah disandarkan. Buku ini sudah berbaloi.

\end{tcbpractice}
